\section{Preserving low pointwise factors and concentrating the remaining defect in degree four}
\label{sec:136-low-factor-preservation}

\paragraph{Archive and scope.}
The uniform hypotheses are cote136, physical pp.~117 and119,
footers116 and118, scans117 and119; the common-source question
is physical p.~121/footer120/scan121. This is a constructive
continuation of ARCH-136-FIRST-CELL-RESIDUES. It controls low
primitive compatibility during actual vertical attachments,
except in form degree four. The resulting degree-four
concentration is a complementary partial source; it is not
asserted to settle the integral common-source question or to
be a logically stronger substitute for the earlier degree-two
and degree-three concentrations.

We use the literal common first skeleton \(F\), the supplied
RH1 construction, and ARCH-136-SKELETON-TATE-LIFTS. In the
latter, degree-two primitive pairs can be chosen with common
graph restrictions; in degree at least four the same is true.
Degree three has the graph-difference obstruction. All cutoffs
in the following argument are fixed from the original integral
targets, so they also work after localization at any prime.

\subsection{The pointwise property}

Write \(S=\Gamma_{\mathbb Z}(T)\). Let \(A\to C_j\)
be finitely many strict maps with first skeleton literally \(F\)
and the identity maps on \(F\). Suppose \(A_0=S\), weight
zero is \(\mathbb Z[0]\), and the rational vertical
cohomology is canonically the coefficient algebra in every
simplex and weight. We impose the following additional property
at each prime \(\ell\) and simplex \(r\), writing
\(S_\ell=S\otimes\mathbb Z_{(\ell)}\).

The algebra \(A_{r,(\ell)}\) has a pointwise grading by
an abelian group, with residue-zero differential PD subalgebra
\(B\). All zeroforms belong to \(B\). Every other
residue component is a contractible differential \(B\)-module,
with a graded \(B\)-linear contraction, and has no zeroforms.
The grading need not be simplicial. The restrictions of the
target maps to \(B\) factor as
\[
 B\xrightarrow{j}D\xrightarrow{\phi_j}(C_j)_{r,(\ell)},
 \qquad \phi_j|_{G_r}\text{ is the same map }D\to F_{(\ell)}(G_r),
 \quad G_r=\operatorname{sk}_1\Delta[r].                 \tag{136.LF1}
\]
Here \(D\) is a graded-free strict PD algebra over \(S_\ell\)
with a special differential, and
\[
                     H^0(D)=S_\ell,\qquad
                     H^1(D)=H^2(D)=H^3(D)=0              \tag{136.LF2}
\]
in every weight. In particular \(D^0\) is integer-torsion-free.
No simplicial structure on the collection of factors \(D\)
is assumed. The trivial residue grading is allowed, in which
case the entire source simplex factors through \(D\).

For clarity, this property implies that every sufficiently
high-weight closed degree-two Moore pair of \(A_r\) has
target primitives with common graph values. Locally, contract
its nonzero residue components in the source, and factor the
remaining component through (136.LF2). This gives common local
graph primitives for every modulus. Integrally, choose arbitrary
target primitives in the fixed comparison range. Their canonical
graph differences lie in \(H^1(G_r;S^w/NS^w)\), vanish
at every prime, and therefore vanish over \(\mathbb Z\).
Graph cone primitives extend by the fixed acyclic-row bounds
and make the restrictions literally equal. This is exactly the
local-to-integral argument in FIRST-CELL-RESIDUES.

\paragraph{ARCH-136-LOW-FACTOR-PRESERVATION.}
The property (136.LF1)--(136.LF2) is preserved by the following
actual strict operations under \(A\), with the retained
target maps and literal first skeleton \(F\): free horizontal
unit-disk cells of simplex dimension at least two; RH1 steps;
and free Tate-pair attachments in form degree
\[
                              q=2,3\quad\text{or}\quad q\geq5.
                                                               \tag{136.LF3}
\]
For \(q=2,3\), the new generator weights must exceed one
fixed low target bound. The Tate attachments are allowed any
already prescribed target maps with common first-skeleton values,
provided their pointwise presentation is free on the new pairs.
In particular graph-relative representable attachments are allowed.
No compatibility conclusion for degree-four attachments is included.

\begin{proof}
We first record that tensoring any factor \(D\) satisfying
(136.LF2) with free unit disks preserves those equations, even
when the old \(D\) is not itself a tensor of unit disks.
The nonconstant blocks of an even lower-degree disk are
contractible. For an odd lower degree, remove its arity-one
unit block. Every remaining block is
\(S_\ell\xrightarrow{s}S_\ell\), starting in form degree
at least three. In the only case that can contribute in degree
three, a cycle of its tensor with \(D\) would be a closed
\(x\in D^0\) with \(sx=0\). But \(Z^0(D)=S_\ell\)
has no integer torsion. Thus there is no such cycle. Products
of two remaining blocks start in degree six. This proves the
claim, with arbitrary ranks by finite support. Graded freeness
and torsion-freeness of degree zero are also retained.

For horizontal attachments, use the degreewise vertical-disk
splitting of the actual horizontal pushout, and tensor \(D\)
with the same independent disks. Their actual target images
have common graph restrictions because the source still has
first skeleton \(F\). Give all new disks residue zero.
The other residue modules undergo base change and keep their
contractions and absence of zeroforms. This argument is pointwise;
no compatibility of the residue grading with faces is used.

For RH1, every one-cycle belongs to the zero-residue part:
the other parts are contractible and have no zeroforms. All
new zeroforms and all identifying zero-form PD relations may
therefore be placed in that part. The other components again
undergo actual module base change; the scalar PD identity shows
that the full PD ideal is the module ideal generated from the
zero-residue relations. The contractions extend with the usual
differential cancellation. In each repaired weight the exact
replacement sequence is
\[
           0\longrightarrow\underline{S_\ell^w}
             \longrightarrow V^w\xrightarrow{d_V}Z^1(A^{w,*}_{(\ell)})
             \longrightarrow0.                          \tag{136.LF4}
\]
This retains the actual, possibly nonzero, normalized boundary
\(V_1\to S_\ell\). Choose a vertex. The common graph
evaluation in (136.LF1) gives an augmentation \(\eta:D\to S_\ell\).
Map a new \(x\in V_r\) to the unique zeroform \(y\in D\)
with \(dy=j(d_Vx)\) and \(\eta(y)\) equal to the
prescribed vertex value of \(x\). Equations (136.LF2) give
existence and uniqueness. This agrees with old zeroforms, respects
the PD pushout, and agrees with every actual retained target
value by equality of differential and vertex value in the fixed
target \(H^0\) range. This is the full RH1 factor argument,
without assuming that (136.LF4) splits simplicially.

It remains to treat Tate pairs. Work at one prime and simplex,
and write one such pointwise pair as
\[
 |c|=q-1,\quad |t|=q-2,\qquad dc=z,\quad dt=Nc-b,
                    dz=0,\quad db=Nz.                    \tag{136.LF5}
\]
The old pair \((b,z)\) is a closed Moore-cone element of
degree \(q-1\). Decompose it by the old residue grading.
Every nonzero component has an actual cone primitive from
its module contraction: the cone homotopy is
\((x,y)\mapsto(hx,-hy)\). Thus there are old elements
\(c_\chi,t_\chi\) with
\(dc_\chi=z_\chi\), \(dt_\chi=Nc_\chi-b_\chi\).
Replace the new coordinates by
\[
                  \bar c=c-\sum_{\chi\ne0}c_\chi,\qquad
                  \bar t=t-\sum_{\chi\ne0}t_\chi.       \tag{136.LF6}
\]
Only finitely many components occur. Translation by old elements
is an invertible free strict PD coordinate change, with the
specified odd exterior rule. The new equations are (136.LF5)
with \(z,b\) replaced by their zero-residue parts. The
target values of the translated generators still have common
graph restrictions, since all subtracted elements are old source
elements. Assign the translated generators residue zero. The
new zero part is the corresponding free Tate extension of \(B\),
and every other component is its module base change. Its old
contraction persists. Hence we only have to extend the factor
of this zero part.

For \(q=2,3\), choose in \(D\)
\[
              dc^*=j(z),\qquad dt^*=Nc^*-j(b),            \tag{136.LF7}
\]
using (136.LF2). For the prescribed target values define
\[
 \delta c_j=f_j(\bar c)-\phi_j(c^*),\qquad
 \delta t_j=f_j(\bar t)-\phi_j(t^*).
                                                               \tag{136.LF8}
\]
They satisfy \(d\delta c_j=0\),
\(d\delta t_j=N\delta c_j\), and have common graph
values. Their weights are the weights of the new generators.

When \(q=2\), choose zeroforms \(\alpha_j\) with
\(d\alpha_j=\delta c_j\). Differences of their graph
restrictions are closed zeroforms, hence common coefficients
on the connected graph. Adjust by those coefficients to make
all graph restrictions of \(\alpha_j\) equal. Then
\(\lambda_j=\delta t_j-N\alpha_j\) is a target coefficient,
and its common graph value implies that all \(\lambda_j\)
are the same coefficient \(\lambda\). Adjoin one free unit
disk \(du=v\) in forms zero and one to \(D\), evaluated
by \(u\mapsto\alpha_j\). Extend the factor by
\[
                     \bar c\longmapsto c^*+v,\qquad
                     \bar t\longmapsto t^*+Nu+\lambda.
                                                               \tag{136.LF9}
\]
These formulas satisfy the differential equations and give
exactly the prescribed target values. All new factor evaluations
have common graph restrictions.

When \(q=3\), use the degree-two common-skeleton primitive
lemma on (136.LF8). It gives \(\alpha_j\) in form one and
\(\beta_j\) in form zero, with common graph restrictions, such that
\[
                d\alpha_j=\delta c_j,\qquad
                d\beta_j=N\alpha_j-\delta t_j.           \tag{136.LF10}
\]
Here is the precise coefficient normalization in that use.
Arbitrary cone primitives have degree zero; their graph
differences have cohomology in
\(H^0(P_N(F(G_r)))=S_\ell^w/NS_\ell^w\), not zero.
Subtract a coefficient representative of each difference.
The remaining difference is the differential of a degree-minus-one
cone cochain, which is an ordinary graph zeroform. Extend that
zeroform to the target using its acyclic form-zero row, and
subtract its differential as well. This makes the primitive
restrictions literally equal. The graph identity follows from
\(H^0(F(G_r))=S_\ell\), torsion-freeness of its ordinary
\(H^1\), and the fixed endpoint bound. There is no claim
that the degree-zero cone cohomology vanishes. The negative
cone cohomology is zero because multiplication by \(N\)
is injective on the coefficient group.

Adjoin free unit disks \(du=v\), \(du'=v'\) in forms
\((1,2)\), \((0,1)\), evaluated by
\(u\mapsto\alpha_j\), \(u'\mapsto\beta_j\).
Set
\[
                    \bar c\longmapsto c^*+v,\qquad
                    \bar t\longmapsto t^*+Nu-v'.         \tag{136.LF11}
\]
Again the differential and the prescribed evaluations agree
exactly. The preliminary unit-disk calculation proves (136.LF2)
for the enlarged factor in both cases. These arguments use only
the targets' fixed comparisons through degree two, form-zero
row extension, and the common graph bound. They are uniform
in \(N\), simplex dimension, and the number of pairs.

For \(q\geq5\), do not take a primitive in the old factor.
Adjoin a formal free Tate pair \(c^*,t^*\) to \(D\) with
the equations (136.LF7), and evaluate it at the already prescribed
target images of \(\bar c,\bar t\). These evaluations are
common on the graph. The resulting factor remains graded-free,
and its degree-zero part is unchanged. For \(q\geq6\)
its cohomology through degree three is visibly unchanged. For
\(q=5\), a possible new three-cycle is
\(x+\sum_i f_i t_i^*\), with \(x\in D^3\), \(f_i\in D^0\).
The independent coefficient of \(c_i^*\) in its differential
is \(N_i f_i\). Since \(D^0\) is integer-torsion-free,
every \(f_i=0\). Thus all three-cycles and their boundaries
are old, and lower forms are unchanged. This proves (136.LF2).
Mixed terms of two new generators have form at least six and
do not enter this check.

All constructions apply independently to simultaneous free
pairs. They are maps on the actual pointwise presentations of
the strict simplicial pushouts. Their evaluations are the retained
simplicial maps; the auxiliary factors themselves need not be
simplicial. This proves the assertion at every prime. RH1 and
all free presentations commute with localization, so the local
arguments apply to the actual integral operations, not to newly
chosen incompatible source algebras. Compatible factors and
module contractions extend through filtered colimits by finite
support. This completes the preservation proof.
\end{proof}

\subsection{The constructive consequence}

\paragraph{ARCH-136-DEGREE-FOUR-CONCENTRATION.}
For every finite family of admissible integral targets, after
the supplied common-first-skeleton replacement, there is a strict
common source \(A^\infty\) with first skeleton exactly \(F\),
rational coefficient cohomology in every weight, all uniform
horizontal tails, and fixed vertical bounds satisfying
\[
\begin{gathered}
 H^0((A^\infty)_r^{w,*})=S^w\quad(w\geq\lambda_0),\\
 H^q((A^\infty)_r^{w,*})=0
       \quad(q\geq1, q\ne4, w\geq\lambda_q),\\
 H_s\bigl(N(A^\infty)_*^{w,q}\bigr)=0
       \quad(w\geq\mu_q, s\geq0).
\end{gathered}                                             \tag{136.LF12}
\]
The remaining \(H^4\) groups are integer torsion. The source
has the controlled pointwise factor property (136.LF1)--(136.LF2)
at every prime. More precisely, the same construction can be made
under any initial source satisfying the hypotheses and that property.
In particular it can retain a genuine guarded native first cell
of ARCH-136-FIRST-CELL-RESIDUES and all its prescribed target maps.

\begin{proof}
The initial source \(L(F)\) has the property with the trivial
residue grading and \(D=A_r\) at every prime. The guarded
first-cell source has it by FIRST-CELL-RESIDUES, with its native
residue grading at \(p\) and the trivial grading at other
primes. More generally start from the source in the statement.

Choose the horizontal cutoffs from the targets once, with form
zero high enough that the common edge already fills every
coefficient there. Only cells of simplex dimension at least two
are then needed. Choose the RH1 cutoff once. For every
\(q\geq2\), \(q\ne4\), fix a vertical cutoff \(\lambda_q\)
large enough for the target primitive comparisons and their
required graph-extension rows. For \(q=2,3\), also include
the single low bound in LOW-FACTOR-PRESERVATION. No cutoff
depends on the construction stage, simplex dimension, modulus,
or number of cells.

At one vertical step take every closed \(z\) of degree
\(q\ne4\), \(q\geq2\), in its chosen weight tail.
Rational coefficient cohomology gives an integer \(N>0\)
and \(b\) with \(db=Nz\). Choose target primitive pairs.
For \(q=2\) and \(q\geq5\) they can be given common
graph values by the common-skeleton Tate-lift theorem. For
\(q=3\) the pointwise factor property gives the same conclusion
by the local-to-integral argument following (136.LF2). Those
graph values are actual sections of the old source first skeleton
\(F\). Attach the free Tate pairs on the representable simplices,
identify their graph values with those old elements, and use the
chosen target sections. This is an actual differential-stable
strict PD pushout. Pointwise it is free on the relative basis
simplices, so the preservation theorem applies to its actual
target maps in every degree. The first skeleton stays literally
\(F\). In simplex dimensions zero and one the positive
cohomology already vanishes; those need no new cells.

Each such pair is rationally a free unit-disk extension by the
change \(c'=c-b/N\), \(dt=Nc'\). Horizontal steps have
the same rational preservation, and RH1 is a rational isomorphism.
Thus the rational hypothesis needed at every stage persists.
Alternate horizontal steps, these vertical steps, and RH1, always
with the chosen bounds. LOW-FACTOR-PRESERVATION supplies the
same pointwise property at each stage, with compatible factor
extensions at each prime and simplex.

Every element has finite support. Every old horizontal cycle in
its tail is killed at a subsequent horizontal step; every old
vertical cycle in a degree other than zero, one, and four is
killed at a subsequent vertical step; and RH1 handles degrees
zero and one. Filtered colimits are exact on the underlying
modules and commute with fixed-degree cohomology. Coefficients
remain injective by their face retractions from simplex zero.
This proves (136.LF12). Rationalization commutes with cohomology,
so the remaining positive cohomology consists of integer-torsion
classes. The compatible pointwise factors and their contractions
pass to the colimit by the same finite-support argument.
\end{proof}

\paragraph{The precise degree-four stop.}
For a degree-four Tate pair the new factor generator has form
two, and adjoining the formal pair can create degree-three
cohomology in the factor. If one instead first enlarges the
factor to make its degree-four image exact, the discrepancy
between a chosen factor primitive and the actual prescribed
simplicial target images is a closed degree-two Moore cone pair.
Correcting that discrepancy requires the degree-three graph
compatibility condition, which is not automatic for an arbitrary
tuple of target differences. Exactness of the old factor through
degree four alone does not make the pointwise choices simplicially
compatible. This is the step omitted in (136.LF3), and no
unrestricted degree-four preservation or full admissible source
is claimed.

\paragraph{Verification status.}
The author completed a whole saved-proof reading, including the
unit-disk calculation for a general factor, the actual RH1 replacement,
strict coordinate translations of contracted components, both correction
formulas and the nonzero coefficient cone class, the degree-five leading
coefficient check, all fixed cutoffs, graph-relative presentations and
filtered colimits. Independent saved-proof reviews are recorded separately.
The degree-four concentration is a complementary construction, with no
strict main-gate movement claim.
